%=========== ΛΥΣΗ - 3ο ΘΕΜΑ Δ ===========
\begin{thema}{Δ}
Θεωρούμε τη συνάρτηση
\[
h(t)=t^3+3t,\quad t\in\mathbb{R}.
\]
Η $h$ είναι παραγωγίσιμη στο $\mathbb{R}$ και
\[
h'(t)=3t^2+3>0
\]
για κάθε $t\in\mathbb{R}$, οπότε είναι γνησίως αύξουσα και συνεπώς συνάρτηση 1-1.
\begin{erwthma}
\item Έστω $x_1,x_2\in\mathbb{R}$ με $f(x_1)=f(x_2)$. Τότε
\begin{align*}
f^3(x_1)+3f(x_1)&=f^3(x_2)+3f(x_2)\\
\Leftrightarrow x_1+4&=x_2+4\\
\Leftrightarrow x_1&=x_2.
\end{align*}
Άρα η $f$ είναι συνάρτηση 1-1.

Επιπλέον, για κάθε $y\in\mathbb{R}$ θέτουμε
\[
x=y^3+3y-4.
\]
Από τη δοσμένη σχέση έχουμε
\[
h(f(x))=f^3(x)+3f(x)=x+4=y^3+3y=h(y).
\]
Επειδή η $h$ είναι 1-1, προκύπτει $f(x)=y$. Επομένως η $f$ έχει σύνολο τιμών το $\mathbb{R}$ και άρα αντιστρέφεται στο $\mathbb{R}$.

Από τη σχέση
\[
f^3(x)+3f(x)=x+4
\]
θέτουμε $y=f(x)$ και παίρνουμε
\[
x=y^3+3y-4.
\]
Άρα
\[
{f^{-1}(x)=x^3+3x-4,\quad x\in\mathbb{R}}.
\]

\item Η $f^{-1}$ είναι παραγωγίσιμη στο $\mathbb{R}$ και
\[
\left(f^{-1}\right)'(x)=3x^2+3>0
\]
για κάθε $x\in\mathbb{R}$. Επομένως η $f^{-1}$ είναι \textbf{γνησίως αύξουσα} στο $\mathbb{R}$.

Επειδή και η αντίστροφή της $f$ είναι γνησίως αύξουσα, η $f$ είναι επίσης γνησίως αύξουσα. Έστω ότι το σημείο $M(\alpha,\beta)$ είναι κοινό σημείο των $C_f$ και $C_{f^{-1}}$. Τότε
\[
f(\alpha)=\beta
\quad\text{και}\quad
f^{-1}(\alpha)=\beta\Leftrightarrow f(\beta)=\alpha.
\]
Αν ήταν $\alpha<\beta$, επειδή η $f$ είναι γνησίως αύξουσα, θα είχαμε
\[
f(\alpha)<f(\beta)\Leftrightarrow \beta<\alpha,
\]
άτοπο. Ομοίως αποκλείεται η περίπτωση $\alpha>\beta$. Άρα $\alpha=\beta$, δηλαδή κάθε κοινό σημείο βρίσκεται στην ευθεία $y=x$.

Επομένως αναζητούμε τις λύσεις της εξίσωσης
\begin{align*}
f^{-1}(x)&=x\\
\Leftrightarrow x^3+3x-4&=x\\
\Leftrightarrow x^3+2x-4&=0.
\end{align*}
Θεωρούμε τη συνάρτηση
\[
g(x)=x^3+2x-4.
\]
Είναι γνησίως αύξουσα στο $\mathbb{R}$, αφού
\[
g'(x)=3x^2+2>0.
\]
Επίσης
\[
g(1)=-1<0
\quad\text{και}\quad
g(2)=8>0.
\]
Άρα, από το Θεώρημα Ενδιάμεσων Τιμών και τη γνήσια μονοτονία της $g$, η εξίσωση $g(x)=0$ έχει ακριβώς μία ρίζα $\alpha\in(1,2)$. Συνεπώς οι $C_f$ και $C_{f^{-1}}$ έχουν μοναδικό κοινό σημείο το
\[
{M(\alpha,\alpha)},
\]
όπου $\alpha$ είναι η μοναδική ρίζα της εξίσωσης $x^3+2x-4=0$.

\item Από τη σχέση
\[
f^3(x)+3f(x)=x+4
\]
προκύπτει ότι
\[
f(x)\to+\infty.
\]
Πράγματι, για κάθε $A>0$, αν
\[
x>A^3+3A-4,
\]
τότε
\[
h(f(x))=x+4>A^3+3A=h(A).
\]
Επειδή η $h$ είναι γνησίως αύξουσα, έχουμε $f(x)>A$.

Επίσης
\begin{align*}
\frac{f(x)}{x}
&=\frac{f(x)}{f^3(x)+3f(x)-4}\\
&=\frac{1}{f^2(x)+3-\dfrac{4}{f(x)}}\longrightarrow 0.
\end{align*}
Διαιρώντας τη δοσμένη σχέση με $x$, έχουμε
\[
\frac{f^3(x)}{x}+3\frac{f(x)}{x}=1+\frac{4}{x}.
\]
Άρα
\[
\left(\frac{f(x)}{\sqrt[3]{x}}\right)^3
=\frac{f^3(x)}{x}
=1+\frac{4}{x}-3\frac{f(x)}{x}\longrightarrow 1.
\]
Επειδή $f(x)>0$ για κάθε αρκετά μεγάλο $x$, συμπεραίνουμε ότι
\[
{\lim_{x\to+\infty}\frac{f(x)}{\sqrt[3]{x}}=1}.
\]

\item Θέτουμε
\[
u=f(x).
\]
Τότε, από τον τύπο της αντίστροφης συνάρτησης,
\[
x=f^{-1}(u)=u^3+3u-4
\]
και
\[
\d x=(3u^2+3)\d u.
\]
Για τα άκρα ολοκλήρωσης έχουμε
\[
f(-4)=0,
\]
διότι $0^3+3\cdot0=-4+4$, και
\[
f(0)=1,
\]
διότι $1^3+3\cdot1=0+4$.
Επομένως
\begin{align*}
\int_{-4}^{0}f(x)\d x
&=\int_0^1 u(3u^2+3)\d u\\
&=\int_0^1(3u^3+3u)\d u\\
&=\left[\frac{3u^4}{4}+\frac{3u^2}{2}\right]_0^1\\
&=\frac{3}{4}+\frac{3}{2}
={\frac{9}{4}}.
\end{align*}
\end{erwthma}
\end{thema}
%=========== ΤΕΛΟΣ ΛΥΣΗΣ - 3ο ΘΕΜΑ Δ ===========
